\documentclass[14pt]{exam}

\usepackage{amsmath}
\usepackage{nicefrac}
\usepackage{amssymb}
\usepackage{xcolor}
\usepackage{diagbox}
\usepackage{float}

\title{Probability and Statistics. Week 1}
\date{}

\def\Var{{\textrm{Var}}\,}
\def\E{{\textrm{E}}\,}
\def\Exp{{\textrm{Exp}}\,}

\begin{document}
	\maketitle
	
	
	
	\begin{questions}
		
		\question
		(Walpole 1.2)
		%According to the journal \textit{Chemical Engineering}, an important property of a fiber is its water absorbency. A random sample of 20 pieces of cotton fiber was taken and the absorbency on each piece was measured. The following are the absorbency values:
		The reaction time, in milliseconds, of 20 amateur racers is measured. A sample is as follows:
		

		$$
		\begin{tabular}{lllllll}
			187.1 & 214.1 & 207.2 & 218.1 & 192.9 & 224.3 & 201.7\\
			237.1 & 194.4 & 205.0 & 189.2 & 203.3 & 230.0 & 228.5\\
			192.5 & 217.7 & 221.1 & 197.7 & 180.4 & 211.2
		\end{tabular}
		\label{tab:1_2}
		$$
		
		\begin{parts}
			\part Calculate the sample mean and sample median for the above sample values.
			\part Compute the 10\% trimmed mean.
			\part Do a dot plot of the reaction time data.
			\part Using only the values of the mean, median, and trimmed mean, do you have evidence of outliers in the data?
		\end{parts}
		
		
		\question
		\label{ex:1_5}
		(Walpole 1.5 and 1.11)
		%Twenty adult males between the ages of 30 and 40 participated in a study to evaluate the effect of specific health regimen involving diet and exercise on the blood cholesterol. Ten were randomly selected to be in a control group, and ten others were assigned to take part in the regimen as the treatment group for a period of 6 months. The following data show the reduction in cholesterol experienced for the time period for 20 subjects:
		The effect of regular exercising on a cholesterol level is being evaluated. All participants are  divided into two equal groups. Reductions in cholesterol levels reported below.
		$$
		\begin{tabular}{cccccc}
			Control group: & 7 & 3 & -4 & 14 & 2\\
						   & 5 & 22 & -7 & 9 & 5\\
			Treatment group: & -6 & 5 & 9 & 4 & 4\\
							 & 12 & 37 & 5 & 3 & 3
		\end{tabular}
		\label{tab:1_5}
		$$
		\begin{parts}
			\part Do a dot plot of the data for both groups on the same graph.
			\part Compute the mean, median, and 10\% trimmed mean for both groups.
			\part Explain why the difference in means suggests one conclusion about the effect of the regimen, while the difference in medians or trimmed means suggests a different conclusion.
			\part Compute the sample variance and the sample standard deviation for both groups.
		\end{parts}
		
		\question
		(Walpole 1.13)
		A manufacturer of electronic components is interested in determining the lifetime of a certain type of battery. A sample, in hours of life, is as follows:
		
		$$
		123, 116, 122, 110, 175, 126, 125, 111, 118, 117
		$$
		
		\begin{parts}
			\part Find the sample mean and median.
			\part What feature in this data set is responsible for the substantial difference between the two? Recalculate the mean without that feature and compare the results.
		\end{parts}
		
		\question
		(Walpole 1.17)
		A study of the effects of smoking on sleep patterns is conducted. The measure observed is the time, in munutes, that it takes to fall asleep. These data are obtained:
		
		$$
		\begin{tabular}{ccccc}
			Smokers (A): & 69.3 & 56.0 & 22.1 & 47.6\\
					& 53.2 & 48.1 & 52.7 & 34.4\\
					& 60.2 & 43.8 & 23.2 & 13.8\\
			Nonsmokers (B):	& 28.6 & 25.1 & 26.4 & 34.9\\
						& 29.8 & 28.4 & 38.5 & 30.2\\
						& 30.6 & 31.8 & 41.6 & 21.1\\
						& 36.0 & 37.9 & 13.9
		\end{tabular}
		\label{tab:1_17}
		$$
		
		\begin{parts}
			\part Find the sample mean for each group.
			\part Find the sample standard deviation for each group.
			\part Make a dot plot of the data sets A and B on the same line.
			\part Comment on what kind of impact smoking appears to have on the time required to fall asleep.
		\end{parts}
		
		\question
		(Walpole 1.18) % b & c only
		The following scores represent the final examination grades for an elementary statistics course:
		
		$$
		\begin{tabular}{ccccccccc}
			23 & 60 & 79 & 32 & 57 & 74 & 52 & 70 & 82\\
			36 & 80 & 77 & 81 & 95 & 41 & 65 & 92 & 85\\
			55 & 76 & 52 & 10 & 64 & 75 & 78 & 25 & 80\\
			98 & 81 & 67 & 41 & 71 & 83 & 54 & 64 & 72\\
			88 & 62 & 74 & 43 & 60 & 78 & 89 & 76 & 84\\
			48 & 84 & 90 & 15 & 79 & 34 & 67 & 17 & 82\\
			69 & 74 & 63 & 80 & 85 & 61
		\end{tabular}
		\label{tab:1_18}
		$$
		
		\begin{parts}
			\setcounter{partno}{1}
			\part Construct a relative frequency histogram, and discuss the skewness of the distribution.
			\part Compute the sample mean, sample median, and sample standard deviation.
		\end{parts}
		
		\question
		(Walpole 1.20) % b, c, d only
		The following data represents the length of life, in seconds, of 50 fruit flies subject to a new spray in a controlled laboratory experiment:
		
		$$
		\begin{tabular}{cccccccccc}
			17 & 20 & 10 & 9 & 23 & 13 & 12 & 19 & 18 & 24\\
			12 & 14 & 6 & 9 & 13 & 6 & 7 & 10 & 13 & 7\\
			16 & 18 & 8 & 13 & 3 & 32 & 9 & 7 & 10 & 11\\
			13 & 7 & 18 & 7 & 10 & 4 & 27 & 19 & 16 & 8\\
			7 & 10 & 5 & 14 & 15 & 10 & 9 & 6 & 7 & 15
		\end{tabular}
		\label{tab:1_20}
		$$
		
		\begin{parts}
			\setcounter{partno}{1}
			\part Set up a relative frequency distribution.
			\part Construct a relative frequency histogram.
			\part Find the median.
		\end{parts}
		
		\question
		Implement in code:
		\begin{itemize}
			\item Mean computation;
			\item Median computation;
			\item Standard deviation;
			\item MAD computation;
			\item Histogram plotting.
		\end{itemize}
		
		You can download required ipynb notebook and data sample from Moodle.
		
		\textit{Note:} you can use Google Colab or configure local Jupyter notebook environment.
		
		\question
		Rolling a fair die you have equal chances to observer 1, 2, 3, 4, 5, or 6.
		
		In case you roll 2 dice, resulting sum is between 2 and 12. Both 9 and 10 can be obtained in two ways: $9 = 3 + 6 = 4 + 5$, $10 = 4 + 6 = 5 + 5$.
		
		In case you roll 3 dice, both 9 and 10 obtained in 6 ways. 
		
		Then, why 9 shows is more likely to be observed when you roll 2 dice, and 10 is more frequently observed when you roll 3 dice?
		
		\textit{Note:} word "die" is a singular form of "dice". However, nowadays those "dice" is used both as a singular and plural.
		
		\question
		When you roll a die 4 times probability that 1 shows up at least once exceeds $\nicefrac{1}{2}$. However, when you roll 2 dice 24 times, probability that 1 appears on both dice simultaneously at least once is below $\nicefrac{1}{2}$. This result is contradictory, because probability of observing 1 is 6 times more than probability of observing two 1's, and 24 rolls is exactly 6 times more than 4 rolls. Explain this contradiction.
		
		\question
		We flip 3 fair coins. At least two of them show the same side, and head and tails on a 3rd coin are equally likely to appear. It follows that all three show the same side is $\nicefrac{1}{2}$. Is it true?
	\end{questions}
\end{document}